\documentclass[11pt,a4paper]{article}
\usepackage[T1]{fontenc}
\usepackage[utf8]{inputenc}
\usepackage[ngerman]{babel}
\usepackage{lmodern}
\usepackage{amsmath,amssymb,amsthm}
\usepackage{mathtools}
\usepackage{enumitem}
\usepackage[a4paper,margin=2.6cm]{geometry}
\usepackage[hidelinks]{hyperref}

\newtheorem{satz}{Satz}
\newtheorem{korollar}[satz]{Korollar}
\newtheorem{problem}[satz]{Problem}
\theoremstyle{remark}
\newtheorem{bemerkung}[satz]{Bemerkung}
\theoremstyle{definition}
\newtheorem{beispiel}[satz]{Beispiel}

\newcommand{\C}{\mathbb{C}}
\newcommand{\R}{\mathbb{R}}
\newcommand{\Z}{\mathbb{Z}}
\newcommand{\N}{\mathbb{N}}
\newcommand{\T}{\mathbb{T}}
\newcommand{\D}{\mathbb{D}}
\newcommand{\TT}{\mathcal{T}}
\newcommand{\eu}{\mathrm{e}}
\newcommand{\im}{\mathrm{i}}
\newcommand{\sperb}{\sigma_{\mathrm{per}}}

\title{Zur Zyklizität: Dobricks Satz und positive Halbgruppen}
\author{Claude (Anthropic)\thanks{Zusammengestellt auf Bitte von Ulrich Groh (Tübingen).}}
\date{16.~Oktober 2026}

\begin{document}
\maketitle

\begin{abstract}
In dem Preprint \cite{dobrick} beweist A.~Dobrick die Zyklizitätsvermutung für einzelne positive Operatoren auf komplexen Banachverbänden.
Diese Notiz geht der Frage nach, ob sich das Resultat oder die Methode auf stark stetige positive Halbgruppen übertragen lässt, also auf die Frage, ob die Randspektren $\sigma_b(A)$ ihrer Generatoren imaginär additiv zyklisch sind.
Diese Frage ist, wie in den Updated Notes zu Part-II und Part-III von LNM~1184 \cite{ln} festgehalten, auch für $C(K)$ offen.
Ergebnis: Für die einzelnen Operatoren $T(t)$ folgt die Zyklizität des peripheren Spektrums sofort; für den Generator nicht.
Die Methode lebt auf dem Kreis $r(T)\T$, und der Weg zum Generator scheitert am Spektralabbildungssatz.
Für $C(K)$ zeigen wir an einem Beispiel, dass ein strikt positives invariantes Maß als Hilfsmittel nicht ausreicht, und geben einen kurzen, elementaren Beweis der (bekannten) Zyklizität im kontraktiven Fall.
\end{abstract}

\section{Ausgangspunkt}

Für $T \in \mathcal{L}(E)$ sei $\sperb(T) \coloneqq \sigma(T) \cap r(T)\T$ mit $\T = \{z \in \C \colon |z| = 1\}$ und $\D = \{z \in \C \colon |z| < 1\}$.

\begin{satz}[{\cite[Theorem~1.2]{dobrick}}]\label{satz:dobrick}
Ist $T$ ein positiver Operator auf einem komplexen Banachverband $E$ mit $r(T) = 1$, so gilt $\lambda^k \in \sigma(T)$ für alle $\lambda \in \sperb(T)$ und alle $k \in \Z$.
\end{satz}

Der Beweis beruht auf drei Bausteinen:
\begin{enumerate}[label=(\Alph*)]
\item \emph{Vererbung unter Domination} \cite[Theorem~2.3]{dobrick}: Aus $0 \le S \le T$ und $r(T) = 1$ folgt
\[
  \sigma\bigl(S + \mu(T-S)\bigr) \cap \T = \sigma(S) \cap \T \subseteq \sigma(T) \cap \T \qquad (\mu \in \D).
\]
Das Argument benutzt den Satz von Vesentini (Subharmonizität von $\log r(F(\mu))$ für holomorphe $F$) und das Maximumprinzip.

\item \emph{Asymptotisch rotations-selbstähnliche Minoranten} \cite[Proposition~3.4]{dobrick}: Ist $Tz = \lambda z$ mit $z \neq 0$, $|\lambda| = 1$, und $U$ ein Torsionsoperator mit $U|z| = z$, so liefert die Iteration von
\[
  \Phi(Q) \coloneqq \tfrac{1}{4}\bigl(2Q + \bar\lambda U^{-1}QU + \lambda UQU^{-1}\bigr)
\]
positive Operatoren $S_n = \Phi^n(T) \le T$ mit $r(S_n) = 1$ und $\|S_nU - \lambda US_n\| \to 0$.

\item \emph{Ultrapotenz und Bidual}: Sie machen aus einer approximativen eine exakte Rotations-Selbstähnlichkeit bzw.\ aus einem approximativen Eigenwert einen Eigenwert, und sie liefern Dedekind-Vollständigkeit (für die Konstruktion des Torsionsoperators über den Signum-Operator).
\end{enumerate}

\section{Was für Halbgruppen sofort folgt}

\begin{korollar}\label{kor:T(t)}
Sei $\TT = (T(t))_{t \ge 0}$ eine positive Halbgruppe auf einem komplexen Banachverband $E$ und $t > 0$ mit $\rho_t \coloneqq r(T(t)) > 0$.
Dann gilt für jedes $\lambda \in \sperb(T(t))$
\[
  \rho_t \Bigl(\frac{\lambda}{\rho_t}\Bigr)^{k} \in \sigma(T(t)) \qquad (k \in \Z),
\]
d.\,h.\ das periphere Spektrum jedes einzelnen Operators $T(t)$ ist zyklisch.
\end{korollar}

\begin{proof}
Satz~\ref{satz:dobrick}, angewandt auf den positiven Operator $T(t)/\rho_t$.
\end{proof}

\section{Warum daraus nichts für \texorpdfstring{$\sigma_b(A)$}{sigma\_b(A)} folgt}

Sei $A$ der Generator von $\TT$.
Wir nehmen $s(A) = \omega_0 = 0$ an.
Für positive Halbgruppen auf $C_0(X)$ ist $s(A) = \omega_0$ immer erfüllt (Part-II, Theorem~3.1.1 in \cite{ln}), auf $L^p$ nach dem Satz von Weis \cite{weis}; auf allgemeinen Banachverbänden dagegen nicht (Beispiel von Arendt auf $L^p \cap L^q$, vgl.\ \cite{abhn}).

Ist $\im\beta \in \sigma_b(A)$, so folgt aus dem Spektralinklusionssatz $\eu^{\im\beta t} \in \sperb(T(t))$, und Korollar~\ref{kor:T(t)} ergibt
\[
  \eu^{\im k\beta t} \in \sigma(T(t)) \qquad (t > 0,\ k \in \Z).
\]
Gewünscht ist aber $\im k\beta \in \sigma(A)$.
Nach Greiner (Part-I, Proposition~3.7.8 in \cite{ln}, nach Reskalierung) gilt für $\mu \in \C$ und $t > 0$:
\[
  \eu^{\mu t} \in \rho(T(t)) \iff
  \begin{cases}
    \mu + \tfrac{2\pi\im}{t}\Z \subseteq \rho(A) \text{ und} \\[2pt]
    \sum_{n \in \Z} R\bigl(\mu + \tfrac{2\pi\im n}{t}, A\bigr)f \text{ ist Cesàro-summierbar für alle } f \in E.
  \end{cases}
\]
Aus $\eu^{\im k\beta t} \in \sigma(T(t))$ folgt also nur: Entweder liegt \emph{irgendein} Punkt von $\im k\beta + \frac{2\pi\im}{t}\Z$ in $\sigma(A)$, oder die Cesàro-Bedingung ist verletzt.
Dass dies für alle $t > 0$ gilt, erzwingt $\im k\beta \in \sigma(A)$ nicht, und zwar schon aus rein mengentheoretischen Gründen:

\begin{beispiel}\label{bsp:menge}
Sei $\beta > 0$, $0 < \delta < \beta$ und $M \coloneqq \im\R \setminus \im(2\beta - \delta, 2\beta + \delta)$.
Dann ist $M$ abgeschlossen, $0, \im\beta \in M$, aber $2\im\beta \notin M$; trotzdem gilt $\eu^{tM} = \T$ für jedes $t > 0$.
Selbst ein exakter Spektralabbildungssatz am Rand, $\sigma(T(t)) \cap \T = \overline{\eu^{t\sigma_b(A)}}$, verträgt sich also mit zyklischen peripheren Spektren aller $T(t)$, ohne dass $\sigma_b(A)$ zyklisch sein müsste.
\end{beispiel}

\begin{bemerkung}
Die Reduktion funktioniert nur, wenn zusätzlich $\sigma_b(A)$ beschränkt ist und $\sigma(T(t)) \cap \T = \eu^{t\sigma_b(A)}$ gilt, etwa bei schließlich normstetigen Halbgruppen.
Ist dann $\sigma_b(A) \subseteq \im[-R,R]$ und $t > 0$ so klein, dass $2\pi/t > R + |k\beta|$, so erzwingt $\eu^{\im k\beta t} = \eu^{t\gamma_t}$ mit $\gamma_t \in \sigma_b(A)$ die Gleichheit $\gamma_t = \im k\beta$.
Zusammen mit der Beschränktheit ergibt das aber nur das Bekannte: $\sigma_b(A) = \{s(A)\}$.
\end{bemerkung}

\section{Warum sich die Methode nicht direkt überträgt}

Alle drei Bausteine (A)--(C) leben auf dem Kreis $r(T)\T$.

\paragraph{Zu (A).}
Für Generatoren bräuchte man die folgende Aussage, für die wir kein Argument sehen:

\begin{problem}\label{prob:vererbung}
Seien $\mathcal{S} = (S(t))_{t\ge0}$ und $\TT = (T(t))_{t\ge0}$ positive Halbgruppen mit Generatoren $B$ bzw.\ $A$ auf einem Banachverband, $0 \le S(t) \le T(t)$ für alle $t \ge 0$ und $s(B) = s(A)$.
Gilt dann $\sigma_b(B) \subseteq \sigma_b(A)$?
\end{problem}

Der Umweg über die Resolventen hilft nicht: Für reelles $\lambda > s(A)$ ist zwar $0 \le R(\lambda,B) \le R(\lambda,A)$ und $r(R(\lambda,B)) = r(R(\lambda,A)) = (\lambda - s(A))^{-1}$, aber ein Punkt $s(A) + \im\beta$ mit $\beta \neq 0$ wird auf $(\lambda - s(A) - \im\beta)^{-1}$ abgebildet, und dieser Punkt liegt im \emph{Inneren} der Spektralkreisscheibe.
Auf dem Kreis sieht man von $\sigma_b(A)$ nur $s(A)$ selbst.

\paragraph{Zu (B).}
Die Mittelung liefert hier
\[
  S_n(t) = \sum_{k} p_{n,k}\, \eu^{-\im k\beta t}\, U^{-k}T(t)U^{k}
\]
(mit $\lambda = \eu^{\im\beta t}$), und das ist für kein $n \ge 1$ eine Halbgruppe, da die Phase von $t$ abhängt.
Die naheliegende Generatorvariante ist
\[
  B_n \coloneqq \sum_{k} p_{n,k}\, A_k, \qquad A_k \coloneqq U^{-k}AU^{k} - \im k\beta .
\]
Jedes $A_k$ erzeugt $\eu^{-\im k\beta t}U^{-k}T(t)U^k$, und diese Halbgruppen sind durch $\TT$ modulus-dominiert: $|\eu^{-\im k\beta t}U^{-k}T(t)U^k x| \le T(t)|x|$.
Mit $\overline{U}\,{=}\,U^{-1}$ (vgl.\ \cite[Proposition~A.3]{dobrick}) ist $\overline{A_k} = A_{-k}$; wegen $p_{n,k} = p_{n,-k}$ ist $B_n$ reell.
Falls der Abschluss von $B_n$ ein Generator ist (die Definitionsbereiche $U^{-k}D(A)$ sind im Allgemeinen verschieden), liefert die Trotter-Formel eine durch $\TT$ modulus-dominierte, reelle Halbgruppe.
Es fehlen jedoch die Positivität (Dobricks Realteil-Argument arbeitet mit dem Operator $\Phi(Q)$, nicht mit Exponentialen von Summen) und die Kritikalität $s(B_n) = s(A)$.

\paragraph{Zu (C).}
Für Halbgruppen muss man sich in der Ultrapotenz auf den Stetigkeitsteil $E^{\mathfrak{U}}_{\TT}$ einschränken (die Halbgruppen-$\mathcal{F}$-Produkte in Part-I, Abschnitt~1.3.7 von \cite{ln}); dort wird $\im\beta$ ein Eigenwert.
Im Bidual geht jedoch die starke Stetigkeit verloren, und gerade dort wird der Torsionsoperator konstruiert.
Auf dem Hauptideal $E_{|z|} \cong C(K)$ existiert der Signum-Operator zwar auch ohne Dedekind-Vollständigkeit (wie in Part-III, Theorem~3.2.2 von \cite{ln}), dort benutzt man aber $T(t)|z| = |z|$, was hier nicht vorliegt.

\section{Der Fall \texorpdfstring{$C(K)$}{C(K)}}

Für $C(K)$ liegt zunächst folgender Ansatz nahe: Man wählt ein strikt positives invariantes Maß $\varphi$ und hofft, approximative Eigenvektoren $f_n$ zu $\im\beta \in \sigma_b(A)$ mit $\inf_n \varphi(|f_n|) > 0$ zu finden.
Im $\mathcal{F}$-Produkt hätte dann der Grenzeigenvektor $h$ positive $\varphi$-Masse, woraus $T_{\mathcal{F}}(t)|h| = |h|$ und mit Part-II, Theorem~3.2.4 die Zyklizität folgen könnte.
Das folgende Beispiel zeigt, dass dies nicht trägt.

\begin{beispiel}\label{bsp:kreise}
Sei $K \coloneqq \{0\} \cup \bigcup_{m \in \N} C_m \subseteq \R^2$ mit den konzentrischen Kreisen $C_m \coloneqq \{x \colon |x| = 1/m\}$; jedes $C_m$ ist offen-abgeschlossen in $K$.
Der Fluss $\phi_t$ drehe $C_m$ mit Winkelgeschwindigkeit $\omega_m \coloneqq 1 + 1/m$ und halte $0$ fest.
Dann ist $T(t)f \coloneqq f \circ \phi_t$ eine stark stetige Markov-Halbgruppe auf $C(K)$ mit $s(A) = 0$, und
\[
  \varphi \coloneqq \sum_{m \in \N} 2^{-m} \lambda_m \qquad (\lambda_m \text{ normiertes Bogenmaß auf } C_m)
\]
ist ein invariantes, strikt positives Wahrscheinlichkeitsmaß (jede Umgebung von $0$ enthält fast alle $C_m$).
\begin{enumerate}[label=(\roman*)]
\item $\im \in \sigma_b(A)$: Für $g_m \coloneqq \eu^{\im\theta}$ auf $C_m$ und $0$ sonst gilt $\|(\im - A)g_m\| = 1/m \to 0$.
\item $\im \notin P\sigma(A)$: Auf $C(C_m)$ ist $\sigma(A) = \im\omega_m\Z \not\ni \im$, da $1/\omega_m = m/(m+1) \notin \Z$; im Fixpunkt $0$ ist $Af(0) = 0$.
\item Jede Folge $(f_n)$ mit $\|f_n\| = 1$ und $(\im - A)f_n \to 0$ erfüllt $\varphi(|f_n|) \to 0$.
Denn $K_M \coloneqq C_1 \cup \dots \cup C_M$ ist offen-abgeschlossen und invariant, auf $C(K_M)$ ist $\im \in \rho(A)$, also $\|\mathbf{1}_{K_M} f_n\| \le c_M \|(\im - A)f_n\| \to 0$ und daher $\limsup_n \varphi(|f_n|) \le \sum_{m > M} 2^{-m}$ für jedes $M$.
\end{enumerate}
Strikt positive Invarianz reicht also nicht: Die relevante Masse wandert in Bereiche, die $\varphi$ kaum sieht.
\end{beispiel}

Verzichtet man dagegen auf ein \emph{festes} invariantes Maß und nimmt zu jedem $f_n$ ein eigenes, nur fast invariantes Maß, so erhält man im kontraktiven Fall einen kurzen, elementaren Beweis der Zyklizität, ohne $\mathcal{F}$-Produkte und ohne Signum-Operatoren.
Für $z \in \C$ sei $z^{[k]} \coloneqq |z|\,(z/|z|)^k$ ($k \in \Z$) und $0^{[k]} \coloneqq 0$.
Die Abbildung $z \mapsto z^{[k]}$ ist Lipschitz-stetig mit einer Konstanten $L_k \le 1 + |k|\pi/2$, und für $|\mu| = 1$ gilt $(\mu z)^{[k]} = \mu^k z^{[k]}$.

\begin{satz}\label{satz:kontraktiv}
Sei $\TT$ eine positive Kontraktionshalbgruppe auf $C(K)$, d.\,h.\ $T(t) \ge 0$ und $\|T(t)\| = \|T(t)\mathbf{1}\| \le 1$, äquivalent $T(t)\mathbf{1} \le \mathbf{1}$, für alle $t \ge 0$.
Dann ist $\sigma(A) \cap \im\R$ imaginär additiv zyklisch.
Insbesondere ist $\sigma_b(A)$ zyklisch, wenn $s(A) = 0$ ist.

Allgemeiner gilt: Ist $s(A) > -\infty$ und die reskalierte Halbgruppe $(\eu^{-s(A)t}T(t))_{t \ge 0}$ kontraktiv, so ist $\sigma_b(A)$ zyklisch.
\end{satz}

\begin{proof}
Sei $\im\beta \in \sigma(A)$ und $k \in \Z$.
Die allgemeinere Aussage folgt durch Anwendung auf die reskalierte Halbgruppe mit Generator $A - s(A)$.
Wegen $\|T(t)\| \le 1$ ist $s(A) \le \omega_0 \le 0$; daher ist $\im\beta$ ein Randpunkt von $\sigma(A)$, also $\im\beta \in A\sigma(A)$.

\emph{Schritt 1 (approximative Eigenvektoren).}
Wähle $f_n \in D(A)$ mit $\|f_n\| = 1$ und $\varepsilon_n \coloneqq \|(A - \im\beta)f_n\| \to 0$, o.\,B.\,d.\,A.\ $\varepsilon_n > 0$.
Wegen $T(t)f - \eu^{\im\beta t}f = \int_0^t \eu^{\im\beta(t-r)}T(r)(A - \im\beta)f \,\mathrm{d}r$ gilt $\|T(t)f_n - \eu^{\im\beta t}f_n\| \le t\varepsilon_n$.

\emph{Schritt 2 (fast invariante Maße).}
Wähle $x_n \in K$ mit $|f_n(x_n)| = 1$, setze $\tau_n \coloneqq \varepsilon_n^{-1/2}$ und
\[
  \psi_n \coloneqq \frac{1}{\tau_n}\int_0^{\tau_n} T(s)'\delta_{x_n} \,\mathrm{d}s \qquad (\text{schwach*-Integral}).
\]
Dann ist $\psi_n \ge 0$, $\psi_n(\mathbf{1}) \le 1$, $\|T(t)'\psi_n - \psi_n\| \le 2t/\tau_n$ und
\[
  \psi_n(|f_n|) = \frac{1}{\tau_n}\int_0^{\tau_n} \bigl(T(s)|f_n|\bigr)(x_n)\,\mathrm{d}s
  \ge \frac{1}{\tau_n}\int_0^{\tau_n} \bigl|(T(s)f_n)(x_n)\bigr|\,\mathrm{d}s
  \ge 1 - \tfrac{1}{2}\tau_n\varepsilon_n \longrightarrow 1 .
\]

\emph{Schritt 3 (duale Kandidaten).}
Sei $g_n \coloneqq f_n^{[k]} \in C(K)$ und $\chi_n \coloneqq \overline{g_n}\,\psi_n \in C(K)'$.
Dann ist $\|\chi_n\| \le 1$ und $\|\chi_n\| \ge \chi_n(g_n) = \psi_n(|f_n|^2) \ge \psi_n(|f_n|)^2 \to 1$.

\emph{Schritt 4 (Kernabschätzung).}
Fixiere $t \ge 0$, schreibe $T = T(t)$, $\mu = \eu^{\im\beta t}$, $(Tf)(x) = \int f\,\mathrm{d}\mu_x$ mit positiven Maßen $\mu_x$, $\mu_x(K) = (T\mathbf{1})(x) \le 1$, und $\delta \coloneqq \|T'\psi_n - \psi_n\|$.
Für $\|f\| \le 1$ ist $\mu^k\chi_n(f) = \mu^k (T'\psi_n)(\overline{g_n}f) + O(\delta)$, also
\begin{align*}
  \bigl|(T'\chi_n - \mu^k\chi_n)(f)\bigr|
  &\le \delta + \int\!\!\int \bigl|(\mu f_n(x))^{[k]} - f_n(y)^{[k]}\bigr| \,\mathrm{d}\mu_x(y)\,\mathrm{d}\psi_n(x) \\
  &\le \delta + L_k \Bigl(\int\!\!\int |\mu f_n(x) - f_n(y)|^2 \,\mathrm{d}\mu_x(y)\,\mathrm{d}\psi_n(x)\Bigr)^{1/2},
\end{align*}
wobei im letzten Schritt die Cauchy--Schwarz-Ungleichung mit $\int\!\!\int \mathrm{d}\mu_x\,\mathrm{d}\psi_n \le 1$ benutzt wurde.
Ausmultiplizieren ergibt mit $Tf_n = \mu f_n + e$, $\|e\| \le t\varepsilon_n$:
\begin{align*}
  \int |\mu f_n(x) - f_n(y)|^2 \,\mathrm{d}\mu_x(y)
  &\le \bigl(T|f_n|^2\bigr)(x) - 2|f_n(x)|^2 + 2t\varepsilon_n + |f_n(x)|^2 \mu_x(K) \\
  &\le \bigl(T|f_n|^2 - |f_n|^2\bigr)(x) + 2t\varepsilon_n .
\end{align*}
Integration gegen $\psi_n$ und $\psi_n(T|f_n|^2 - |f_n|^2) = (T'\psi_n - \psi_n)(|f_n|^2) \le \delta$ liefern
\[
  \|T(t)'\chi_n - \eu^{\im k\beta t}\chi_n\| \le \frac{2t}{\tau_n} + L_k\Bigl(\frac{2t}{\tau_n} + 2t\varepsilon_n\Bigr)^{1/2} \longrightarrow 0,
\]
gleichmäßig für $t$ in kompakten Intervallen.

\emph{Schritt 5 (Schluss).}
Angenommen $\lambda \coloneqq \im k\beta \in \rho(A)$.
Für $\tau > 0$ sei $y_n \coloneqq \int_0^\tau \eu^{\lambda(\tau - s)}T(s)'\chi_n\,\mathrm{d}s$.
Dann ist $y_n \in D(A')$ und $(A' - \lambda)y_n = T(\tau)'\chi_n - \eu^{\lambda\tau}\chi_n \to 0$, während nach Schritt~4 $\|y_n - \tau\eu^{\lambda\tau}\chi_n\| \to 0$, also $\liminf_n \|y_n\| \ge \tau$.
Das widerspricht der Beschränktheit von $R(\lambda, A') = R(\lambda, A)'$.
Also ist $\im k\beta \in \sigma(A)$.
\end{proof}

\begin{bemerkung}
\begin{enumerate}[label=(\alph*)]
\item Satz~\ref{satz:kontraktiv} ist inhaltlich nicht neu: Für beschränkte positive Halbgruppen mit $s(A) = 0$ wächst die Resolvente langsam, und Part-III, Theorem~3.2.10 mit Proposition~3.2.9\,(i) in \cite{ln} liefert die Zyklizität sogar auf beliebigen Banachverbänden.
Neu ist der kurze, rein maßtheoretische Beweis.
\item Die Kontraktivität geht genau zweimal ein: bei der Fast-Invarianz der Cesàro-Mittel (Schritt~2) und bei $\mu_x(K) \le 1$ (Schritt~4).
\item Offen bleibt auf $C(K)$ damit der Fall $s(A) = \omega_0 = 0$ mit unbeschränkter Halbgruppe, in dem die Resolvente nicht langsam wächst und $s(A)$ kein Pol ist.
\end{enumerate}
\end{bemerkung}

\section{Mögliche Ansätze}

\begin{itemize}
\item Der Kern ist Problem~\ref{prob:vererbung}, eine Vererbungsaussage für das Randspektrum von Generatoren.
Gelänge sie, so könnte man versuchen, den Rest über $B_n$ und Trotter-Konstruktionen zu retten.
\item Für $C(K)$ bietet sich eine Variante im Geist von Dobrick an: Man sucht zu einer unbeschränkten positiven Halbgruppe mit $s(A) = 0$ eine \emph{beschränkte} kritische Minorante $0 \le S(t) \le T(t)$ mit $s(B) = 0$.
Für $S$ gilt die Zyklizität nach Satz~\ref{satz:kontraktiv} bzw.\ Part-III, Theorem~3.2.10; es bliebe die Vererbung $\sigma_b(B) \subseteq \sigma_b(A)$, also Problem~\ref{prob:vererbung} in der handlicheren Form mit beschränkter Minorante.
Die Maße $\mu_x$ und $\psi_n$ aus dem Beweis von Satz~\ref{satz:kontraktiv} könnten dabei helfen.
\end{itemize}

\section{Zur Arbeit}

Theorem~2.3 aus \cite{dobrick} ist bestechend einfach und nach unserer Prüfung korrekt, ebenso der Schluss $r(S_n) = 1$ im Beweis von Proposition~3.4 (über $S_1 \le (1-\tfrac1n)T + \tfrac1n S_n$ und Theorem~2.3).
Die Anhänge wurden nicht im Detail nachgerechnet.
Es handelt sich um einen arXiv-Preprint vom 7.~Oktober 2026, noch ohne Begutachtung; laut eigener Erklärung entstand er mit Unterstützung durch KI-Werkzeuge.
Hält das Resultat, sollte es in den Updated Notes zu Part-III von \cite{ln} für den Fall eines einzelnen Operators erwähnt werden; die dortigen Aussagen zur offenen Halbgruppenfrage bleiben richtig.

\begin{thebibliography}{9}
\bibitem{abhn}
W.~Arendt, C.\,J.\,K.~Batty, M.~Hieber, F.~Neubrander,
\emph{Vector-valued Laplace Transforms and Cauchy Problems},
2.~Aufl., Birkhäuser, Basel, 2011.

\bibitem{dobrick}
A.~Dobrick,
\emph{On the cyclicity conjecture},
arXiv:2610.10093v1 [math.FA], 7.~Oktober 2026.

\bibitem{ln}
R.~Nagel (Hrsg.),
\emph{One-parameter Semigroups of Positive Operators},
Lecture Notes in Mathematics 1184, Springer, Berlin, 1986; 2.~Aufl.\ in Vorbereitung.

\bibitem{weis}
L.~Weis,
\emph{The stability of positive semigroups on $L_p$ spaces},
Proc.\ Amer.\ Math.\ Soc.\ \textbf{123} (1995), 3089--3094.
\end{thebibliography}

\end{document}
